\documentclass[10pt,letterpaper]{article}
\usepackage[margin=0.7in]{geometry}
\usepackage[T1]{fontenc}
\usepackage{lmodern,microtype,amsmath,booktabs,array,graphicx,float,xcolor}
\usepackage{fancyhdr,enumitem,hyperref}
\hypersetup{colorlinks=true,linkcolor=blue!45!black,urlcolor=blue!45!black}
\setlength{\parindent}{0pt}\setlength{\parskip}{5pt}\setlength{\headheight}{14pt}
\setlength{\emergencystretch}{4em}
\setlist{nosep,leftmargin=*}
\pagestyle{fancy}\fancyhf{}\lhead{18-344 \quad Lab 2: Memory Hierarchy}\rhead{Maxwell Wu}\cfoot{\thepage}
\newcommand{\fig}[3]{\begin{figure}[H]\centering\includegraphics[width=.92\linewidth]{figures/#1.pdf}\caption{#2}\label{#3}\end{figure}}
\begin{document}
\begin{center}
{\LARGE\bfseries Lab 2: Memory Hierarchy}\par\vspace{5pt}
{\large Set associativity, replacement, and a two-level design search}\par
Maxwell Wu (mwu3)\quad |\quad October 24, 2026
\end{center}
\section*{Overview and evidence status}
The simulator models a parameterized data-cache hierarchy driven by Pin memory
references. The completed Task 4 search evaluates 20 distinct hierarchies on the
SPEC2017 \texttt{602.gcc\_s} and \texttt{641.leela\_s} test workloads, producing 40
benchmark measurements and 18 distinct Destiny cache models. The best measured
hierarchy is an 8 KiB, 64-byte, 4-way LRU L1 and a 512 KiB, 128-byte, 8-way LRU L2.
Its equally weighted mean AMAT is 0.25416 ns, modeled area is 6.26309 mm$^2$, and
estimated total cache power is 152.72 mW at an assumed $10^9$ L1 accesses/s.
This is the best design in the measured set, not a demonstrated global optimum.

\textbf{Submission caveat.} The saved evidence does not fully satisfy Task 2's
three-configuration comparison of all four policies. The corrected collection
contains that comparison at only one L1 geometry. Legacy results cover three
geometries but use tree-PLRU and omit round-robin. This report covers all tasks
while distinguishing those datasets; missing measurements are not invented.

\begin{table}[H]\centering\small
\caption{Evidence available for the handout's requirements.}
\begin{tabular}{p{.10\linewidth}p{.43\linewidth}p{.36\linewidth}}\toprule
Task & Evidence in this report & Qualification \\\midrule
1 & Corrected L1 miss rates at eight geometries; single-level AMAT plots & AMAT derived from L1 counters, not separate single-level runs. \\
2 & Four policies at one corrected geometry; legacy three-geometry plot & Two more corrected geometry comparisons are needed. \\
3 & Parameterized hierarchy and deterministic tests & Demand-only model; no writeback or coherence. \\
4 & 20 designs, Destiny metrics, Pareto and sensitivity plots & Bounded search; one run per benchmark/design. \\\bottomrule
\end{tabular}\end{table}
\subsection*{Workloads and provenance}
The completed session is \texttt{session-3973386ea9622d3cee1b}, under
\texttt{results/task4/}. Its manifest records commands, inputs, source and tool
hashes, and environment. Gcc compiles \texttt{t1.c} with
\texttt{-O3 -finline-limit=50000}; leela runs \texttt{test.sgf}. Each is the sole
command in its selected base/test run directory. Pin instruments the full
process, including library references; there is no instruction-limit sampling
or warmup. Runs use one OpenMP thread, disabled ASLR (successful in this session),
and a deterministic random-replacement seed. Eight designs were reused after
provenance validation and twelve were newly evaluated. Completion was recorded
at 18:57:30 Eastern on October 24, 2026.

\clearpage
\section{Simulator design: Tasks 1--3}
\subsection{Set-associative lookup}
For capacity $C$ bytes, block size $B$, and associativity $A$, the number of sets
is $S=C/(BA)$. Address $a$ maps to block $q=\lfloor a/B\rfloor$, set $q\bmod S$,
and tag $\lfloor q/S\rfloor$. The implementation requires power-of-two block
size, associativity, and set count, with integral capacity division. Each line
stores a valid bit and tag. A lookup scans the selected set and fills an invalid
way before evicting a valid line. A reference spanning blocks generates one
access per touched L1 block. Pin counts a read-modify-write operand once as a
read and once as a write, using predicated callbacks.
\subsection{Replacement policies}
\begin{description}[style=nextline]
\item[Random.] Fixed-seed xorshift32 selects a victim when the set is full. It
ignores reuse, and a single seed does not quantify variability across seeds.
\item[LRU.] A per-set ordered list moves every hit or fill to the most-recent
position and evicts the oldest. Recency tracking costs state and update logic.
\item[Bit-PLRU.] A touch sets one bit per way; when all bits are set, all other
bits are cleared. The first clear bit identifies a victim. This differs from
the tree-PLRU used in the older saved measurements.
\item[Round-robin.] The additional policy uses a per-set pointer that advances
on full-set replacement. Its low-state cyclic selection may avoid random
imbalance but cannot protect hot blocks based on hits.
\end{description}
\subsection{Two-level behavior and boundaries}
Each level has independent capacity, block size, associativity, and policy.
L2 blocks must be a multiple of, and at least as large as, L1 blocks. L1 misses
access L2; L2 misses increment DRAM. Without L2, each L1 miss increments DRAM.
Reads and writes allocate on a miss. L2 demand fetches are modeled as reads,
including fetches caused by L1 store misses. The model has no dirty bits,
writebacks, inclusion invalidations, or writes propagated on L1 hits. It is not
a complete write-through or write-back implementation.

Instruction-cache traffic, physical-address translation, prefetching, coherence,
queueing, overlap, and memory-level parallelism are excluded. Instrumented
virtual addresses drive the model in a single-thread workload setup. AMAT is a
serialized access estimate, not CPU execution time or an official SPEC score.
\subsection{Validation}
\texttt{tests/cache\_test.cpp} checks LRU victims, bit-PLRU behavior distinct from
tree-PLRU, repeatable random replacement, direct-mapped policy equivalence,
cross-block references, write hits, single-level DRAM counts, and larger shared
L2 blocks. These tests passed on ECE before the Pin rebuild. Collection requires
zero process exit status and consistent counters, including $N=H+M$, $N=R+W$,
$N_2=M_1$, and $N_{\mathrm{DRAM}}=M_2$ for two-level runs. It does not perform
SPEC reference-output comparison. Python tests cover parsing, AMAT, budgets,
20-design accounting, output generation, and reuse validation using fixture data.

\clearpage
\section{Task 1: associativity, block size, and AMAT}
With L1 miss rate $m_1$, Destiny hit latency $t_{1h}$ and miss latency $t_{1m}$,
the handout's single-level model is
\begin{equation}
\mathrm{AMAT}_1=(1-m_1)t_{1h}+m_1(t_{1m}+9.5\ \mathrm{ns}).
\end{equation}
L1 decisions in this simulator do not depend on L2 contents or timing. The
corrected L1 counters from the initial eight Task 4 geometries therefore support
a \emph{derived} single-level estimate for the recorded reference streams by
treating every L1 miss as a DRAM access. These are not separate single-level
executions. The plots use this calculation, not the legacy AMAT column.
\fig{task1-amat}{Derived single-level AMAT using corrected L1 counters and Destiny latencies. Lines connect measured points; the 32 KiB curve has no 8-way point.}{f:t1a}
\fig{task1-misses}{Measured L1 miss rates for the same eight corrected geometries.}{f:t1m}
At 8 KiB/64 B, increasing associativity from 1 to 4 lowers gcc's miss rate from
6.202\% to 3.666\% and leela's from 5.114\% to 2.234\%. AMAT falls from 0.7914
to 0.5536 ns for gcc and from 0.6884 to 0.4178 ns for leela. Eight ways reduces
misses further but raises leela's AMAT to 0.4312 ns as hit latency rises. Thus
fewer misses need not mean faster access. At four ways, 64-byte blocks outperform
32-byte blocks for both workloads here. Increasing capacity to 32 KiB with
32-byte blocks also reduces misses, but this limited grid establishes no universal optimum.

\clearpage
\section{Task 2: replacement-policy evaluation}
The intended three configurations are 2 KiB/1-way, 4 KiB/2-way, and 8 KiB/4-way,
all with 64-byte blocks. They provide a direct-mapped control and increasing
capacity/associativity to expose conflict and reuse effects. Since size and ways
change together, compare policies \emph{within} a geometry; cross-geometry
comparisons do not isolate either parameter.
\fig{task2-legacy}{Legacy three-geometry results. Exact commands and benchmark identity were not recorded. The old PLRU label is shown as tree-PLRU; round-robin has no legacy measurement. These are historical observations, not corrected Task 2 completion evidence.}{f:t2old}
At 4 KiB/2-way, legacy LRU and tree-PLRU both report 3.367\% misses versus
random's 4.231\%. At 8 KiB/4-way, rates are 1.880\% (LRU), 1.951\% (tree-PLRU),
and 2.235\% (random). Direct-mapped rates vary from 7.701\% to 7.849\%, although
replacement cannot affect a fixed direct-mapped stream. That variation warns
against causal conclusions from small differences; the records do not identify
its cause.
\fig{task2-corrected}{Corrected four-policy comparison at 8 KiB/64 B/4-way L1. L2 is 512 KiB/64 B/8-way and both levels use the named policy. L1 miss rate is independent of L2 policy in this model.}{f:t2new}
The corresponding two-level mean AMAT is 0.28260 ns (LRU), 0.28322 ns
(bit-PLRU), 0.28595 ns (round-robin), and 0.28852 ns (random). LRU is best here;
bit-PLRU is close. These single-run results do not imply a general ordering.
Changing both policies together does not isolate each level's AMAT contribution.
\textbf{Remaining requirement:} collect all four corrected policies for gcc and
leela at the other two geometries. The legacy plot cannot substitute for this
or be relabeled as bit-PLRU evidence.

\clearpage
\section{Task 4: objective and complete design catalog}
Let $m_2$ be the local L2 miss rate, conditional on an L1 miss. The model is
\begin{align}
\mathrm{AMAT}_2={}&(1-m_1)t_{1h}\\
&+m_1\left[t_{1m}+(1-m_2)t_{2h}+m_2(t_{2m}+9.5\ \mathrm{ns})\right].
\end{align}
DRAM costs $19\times0.5=9.5$ ns. All terms use ns. Ranking minimizes the
arithmetic mean of gcc and leela AMAT with equal benchmark weights; exact ties
use lower area, then lower estimated power. No equivalence tolerance is assumed.

The search starts with eight L1 geometries and a 512 KiB/64 B/8-way LRU L2.
Using the best L1, it adds three matched non-LRU policy pairs, one L2 block-size
alternative, four L2 associativities, and four capacities: $8+3+1+4+4=20$.
The sensitivity curves reuse the baseline. Independent level policies are
supported by the tool, but this budgeted experiment tests matched pairs only.
All designs use at most 4104 KiB of data, below both 5,000,000 bytes and the
script's 5 MiB limit. The winner uses 520 KiB.
\begin{table}[H]\centering\small
\caption{All 20 measured designs. C/B/A/P denotes KiB, block bytes, ways, and policy; L=LRU, R=random, P=bit-PLRU, RR=round-robin. AMAT is in ns, area in mm$^2$, power in mW.}
\begin{tabular}{rllrrr}\toprule
ID & L1: C/B/A/P & L2: C/B/A/P & AMAT & Area & Power \\ \midrule
D01 & 8/32/1/L & 512/64/8/L & 0.2941 & 7.367 & 153.62 \\
D02 & 8/32/4/L & 512/64/8/L & 0.2877 & 7.089 & 128.99 \\
D03 & 8/32/8/L & 512/64/8/L & 0.3099 & 7.093 & 142.95 \\
D04 & 8/64/1/L & 512/64/8/L & 0.2928 & 7.336 & 205.08 \\
D05 & 8/64/4/L & 512/64/8/L & 0.2826 & 7.049 & 148.72 \\
D06 & 8/64/8/L & 512/64/8/L & 0.3047 & 7.041 & 159.10 \\
D07 & 32/32/1/L & 512/64/8/L & 0.3197 & 8.527 & 200.83 \\
D08 & 32/32/4/L & 512/64/8/L & 0.3019 & 7.639 & 167.82 \\
D09 & 8/64/4/R & 512/64/8/R & 0.2885 & 7.049 & 151.68 \\
D10 & 8/64/4/P & 512/64/8/P & 0.2832 & 7.049 & 148.67 \\
D11 & 8/64/4/RR & 512/64/8/RR & 0.2859 & 7.049 & 151.14 \\
D12 & 8/64/4/L & 512/128/8/L & 0.2542 & 6.263 & 152.72 \\
D13 & 8/64/4/L & 512/64/1/L & 0.2923 & 6.784 & 148.03 \\
D14 & 8/64/4/L & 512/64/2/L & 0.2844 & 7.431 & 165.11 \\
D15 & 8/64/4/L & 512/64/4/L & 0.2818 & 6.060 & 145.83 \\
D16 & 8/64/4/L & 512/64/16/L & 0.2834 & 6.483 & 143.13 \\
D17 & 8/64/4/L & 128/64/8/L & 0.2997 & 2.309 & 97.56 \\
D18 & 8/64/4/L & 256/64/8/L & 0.2887 & 3.596 & 113.13 \\
D19 & 8/64/4/L & 1024/64/8/L & 0.2833 & 12.406 & 199.27 \\
D20 & 8/64/4/L & 4096/64/8/L & 0.2941 & 32.203 & 436.70 \\
\bottomrule\end{tabular}

\end{table}
This search emphasizes moderate L1 sizes and reuse of completed measurements
to control runtime. Many capacity/block/policy combinations remain unexplored.
The recommendation is the best measured candidate, not a global optimum.

\clearpage
\section{Pareto tradeoffs and per-benchmark AMAT}
A measured design is Pareto optimal if no other design is no worse in mean AMAT,
area, and estimated total power and strictly better in at least one. Six designs
meet this three-objective criterion. A three-objective Pareto point need not be
on both two-dimensional projected frontiers.
\fig{task4-pareto}{All 20 designs and separate two-dimensional frontiers. Numeric labels identify the six three-objective Pareto designs by Table 2 ID. Lines are visual guides, not interpolated designs.}{f:pareto}
\begin{table}[H]\centering
\caption{Three-objective Pareto set, ordered by mean AMAT.}
\begin{tabular}{lrrr}\toprule
ID & AMAT (ns) & Area (mm$^2$) & Power (mW) \\ \midrule
D12 & 0.2542 & 6.263 & 152.72 \\
D15 & 0.2818 & 6.060 & 145.83 \\
D16 & 0.2834 & 6.483 & 143.13 \\
D02 & 0.2877 & 7.089 & 128.99 \\
D18 & 0.2887 & 3.596 & 113.13 \\
D17 & 0.2997 & 2.309 & 97.56 \\
\bottomrule\end{tabular}

\end{table}
D12 minimizes mean AMAT. Relative to D12, D17 increases AMAT by about 17.9\%
but reduces area by about 63.1\% and power by about 36.1\%. D02 offers another
low-power tradeoff using a different L1 block size. A stricter physical budget
could favor these alternatives; the stated performance-first objective selects D12.

\clearpage
\subsection{Per-benchmark results for the Pareto candidates}
\fig{task4-bars}{Gcc and leela AMAT for the six three-objective Pareto designs. Bars show workload differences, not statistical uncertainty.}{f:bars}
\begin{table}[H]\centering\small
\caption{D12 benchmark-specific performance. DRAM/L1 is the global memory-access fraction, not the local L2 miss rate.}
\begin{tabular}{lrrrr}\toprule
Benchmark & L1 miss (\%) & L2 local miss (\%) & DRAM / L1 (\%) & AMAT (ns) \\ \midrule
gcc & 3.666 & 16.571 & 0.608 & 0.2784 \\
leela & 2.234 & 6.831 & 0.153 & 0.2299 \\
\bottomrule\end{tabular}

\end{table}
For the selected hierarchy, gcc has higher L1 and local L2 miss rates than
leela, and its AMAT is about 21\% higher. The global DRAM fraction is about
0.608\% for gcc versus 0.153\% for leela. These differences show why both
workloads must be measured: an aggregate objective alone can conceal different
memory behaviors.

The reported mean is $(0.2784+0.2299)/2\approx0.2542$ ns, using the unrounded
values in the actual calculation. It gives each benchmark equal weight, rather
than pooling counters and allowing the longer reference stream to dominate.
The bars retain the individual benchmark values for all six Pareto candidates
so alternative application priorities can be assessed directly.

Small global DRAM fractions explain AMAT well below 9.5 ns: DRAM latency is
paid rarely. This does not imply full instructions execute in fractions of a
cycle. The cold-start demand model does not capture overlap or throughput, and
AMAT ratios are not runtime speedups. Without repeated runs or multiple random
seeds, the plotted differences also provide no statistical confidence bounds.

\clearpage
\section{L2 sensitivity and selected hierarchy}
The fixed L1 is 8 KiB/64 B/4-way LRU. Both sweeps use 64-byte L2 blocks and LRU.
The associativity sweep fixes L2 at 512 KiB; the size sweep fixes it at 8 ways.
Their common anchor is D05, not final winner D12 (which has 128-byte L2 blocks).
Each curve changes one architectural parameter, recalculating Destiny costs.
\fig{task4-sensitivity}{Controlled sensitivity experiments. Horizontal axes use base-2 logarithmic spacing with explicit measured ticks.}{f:sens}
At 1, 2, 4, 8, and 16 L2 ways, mean AMAT is 0.29233, 0.28441, 0.28179,
0.28260, and 0.28340 ns. Four ways is best in this sweep, about 3.6\% faster
than one way. Beyond four ways, lookup cost offsets conflict-miss benefits.
At 128, 256, 512, 1024, and 4096 KiB, mean AMAT is 0.29968, 0.28874,
0.28260, 0.28334, and 0.29407 ns. The minimum is 512 KiB. Growing to 4096 KiB
raises AMAT by about 4.1\%, area from 7.049 to 32.203 mm$^2$, and power from
148.72 to 436.70 mW. Filling the entire storage budget is not inherently beneficial.

\subsection{Selected design}
D12 changes the anchor's L2 block size to 128 bytes. Mean AMAT falls from
0.28260 to 0.25416 ns (10.1\%), and area falls from 7.049 to 6.263 mm$^2$.
Power rises from 148.72 to 152.72 mW (2.7\%). Destiny optimizes internal
organizations separately, so area need not vary monotonically with an
architectural parameter. Lower AMAT wins despite slightly higher power.

\clearpage
\section{Implementability: latency, power, and tag overhead}
Destiny uses 2D SRAM at 65 nm, 350 K, the LOP setting, normal access, and
optimization for read latency. These are model estimates, not fabricated-chip
measurements; all configurations share these technology assumptions.
\begin{table}[H]\centering
\caption{Destiny metrics for D12: 8 KiB/64 B/4-way L1; 512 KiB/128 B/8-way L2.}
\begin{tabular}{lrr}\toprule
Metric & L1 & L2 \\ \midrule
Read-hit latency (ns) & 0.2061 & 0.45478 \\
Miss latency (ns) & 0.18265 & 0.24966 \\
Write latency (ns) & 0.12703 & 0.22011 \\
Read-hit energy (nJ) & 0.066654 & 0.80579 \\
Write energy (nJ) & 0.061313 & 0.78358 \\
Leakage (mW) & 1.5043 & 57.814 \\
Area (mm$^2$) & 0.14333 & 6.1198 \\
Tag-array area (mm$^2$) & 0.039901 & 0.78592 \\
\bottomrule\end{tabular}

\end{table}
\textbf{Latency.} L1 read-hit latency is 0.20610 ns, or 0.412 cycles at an
illustrative 2 GHz. Rounding to at least one cycle is compatible with the
few-cycle L1 target. L2 read-hit latency is 0.45478 ns. Core routing, arbitration,
pipeline alignment, and controller delay are not fully represented; sub-cycle
array estimates are not fractional-cycle architectural guarantees. They support
plausibility, not timing closure.

\textbf{Energy and power.} For each level $\ell$, the runner estimates
\begin{equation}
E_\ell=R_{h,\ell}e_{h,\ell}+W_{h,\ell}e_{w,\ell}
       +M_\ell(e_{m,\ell}+e_{w,\ell}),
\end{equation}
including a fill write for each miss. Dividing $E_1+E_2$ by L1 accesses gives
demand energy per access. With energy in nJ and offered rate $f$ in accesses/s,
\begin{equation}
P_{\rm dynamic}[\mathrm{mW}]=\overline E[\mathrm{nJ/access}]f\,10^{-6}.
\end{equation}
At $f=10^9$, 0.093404 nJ/access yields 93.404 mW dynamic power. Adding
59.319 mW leakage gives 152.722 mW. L1/L2 write energies are 0.061313/0.783585
nJ. Sub-nJ operations and about 0.153 W do not indicate obviously excessive
cache-only demand, but no chip-wide budget was supplied. Dynamic power scales
linearly with assumed activity; DRAM, writebacks, coherence, and replacement
logic are excluded. Pin wall-clock time is not a hardware power measurement.

\textbf{Metadata.} For 64-bit addresses, $T=64-\log_2 B-\log_2 S$ tag bits.
L1 has 128 lines, 32 sets, and 53-bit tags; L2 has 4096 lines, 512 sets, and
48-bit tags. With one valid bit per line and estimated LRU state of
$A\log_2A$ bits per set, overhead is 896 bytes (10.94\%) for L1 and 26,624
bytes (5.08\%) for L2: 27,520 bytes total. Dirty and coherence bits are excluded
because they are not modeled. Destiny's tag-area fractions, 27.84\% and 12.84\%,
also reflect peripheral circuits and differ from logical bit overhead. Even
including estimated metadata, D12 uses only 560,000 bytes, below 5 MB.

\clearpage
\section{Reproducibility and limitations}
The repository supplies the simulator, \texttt{script.py}, \texttt{task4.py},
tests, and README. Prepare gcc/leela SPEC test workloads on Linux ECE and enter
their input directories and argument arrays in \texttt{task4\_config.json}:
\begin{verbatim}
python3 task4.py --check
python3 task4.py
\end{verbatim}
The default \texttt{compact20} search is reported here; repeating the command
resumes compatible completed jobs. Runs use private input copies. The committed
export retains counters, commands, completion records, and Destiny inputs and
outputs, but excludes licensed SPEC binaries and input copies. The manifest
retains ECE paths for provenance; a TA must substitute locally available
workloads. The exported run index uses relative paths for its saved statistics.

Rebuild this report's figures without rerunning benchmarks:
\begin{verbatim}
python3 -m pip install matplotlib
python3 report/build_figures.py
cd report
pdflatex main.tex
pdflatex main.tex
\end{verbatim}
Tectonic is an alternative compiler. Derived Task 1 values are saved in
\texttt{generated/task1-derived.csv}. All seven figures are vector PDFs with
PNG previews, redrawn from numerical records rather than HTML screenshots.

\subsection{Limits on conclusions}
There is one run per benchmark/design and one random seed; small differences
lack confidence intervals. Only test inputs are used, so train/ref inputs or
other applications may select a different design. The sweep is not exhaustive,
Destiny separately optimizes organizations, and equal benchmark weighting is
a choice rather than an inherent objective. Simplified writes, virtual addresses,
and omitted overlap limit real-processor claims. Counter consistency and exit
status do not substitute for SPEC output-equivalence validation.

\subsection{Remaining submission work}
Task 4 data and the requested Pareto, per-benchmark, and sensitivity plots are
complete for the declared scope. Task 3 functionality and Task 1 measured miss
rates/derived AMAT are documented. Task 2 still needs corrected four-policy
measurements at two additional geometries on both workloads. Until those are
collected, this is an evidence-qualified report, not a claim that every
experimental requirement is fully complete.

\section*{Sources}
\begin{enumerate}
\item Course handout, \texttt{18344-f25-lab2-memoryhierarchy.pdf}, Tasks 1--4,
pp.~1--3. The document itself is dated September 30, 2024.
\item \texttt{results/task4/session-3973386ea9622d3cee1b/}: manifest, measurement
and hardware tables, design-space table, raw run records, and Destiny outputs.
\item \texttt{results/task1-3/task2/task2.csv}, used only in the labeled
historical policy comparison.
\item Repository simulator and runner source; result-export commit
\texttt{1243851}, compact-search runner commit \texttt{1facf8a}.
\end{enumerate}
\end{document}
